\documentclass[11pt]{article}
\usepackage[margin=1in]{geometry}
\usepackage{amsmath,amssymb,amsthm}
\usepackage[hidelinks]{hyperref}

\newtheorem{theorem}{Theorem}
\newtheorem{lemma}[theorem]{Lemma}
\newtheorem{corollary}[theorem]{Corollary}
\theoremstyle{definition}
\newtheorem{definition}[theorem]{Definition}
\theoremstyle{remark}
\newtheorem{remark}[theorem]{Remark}
\newtheorem{question}{Question}

\newcommand{\F}{\mathbb{F}_2}
\newcommand{\wt}{\operatorname{wt}}
\newcommand{\W}{W}
\newcommand{\Coll}{\mathrm{C}}
\newcommand{\one}{\mathbf{1}}

\title{The exact collision cost of low-weight uniformity\\ for polynomial maps over $\F$}
\author{Daniel Ayomide Shittu\\ {\small Rowan University}}
\date{Preprint, version 2, November 2026}

\begin{document}
\maketitle

\begin{abstract}
Let $Q\colon\F^N\to\F^N$ have algebraic degree at most $k$ in each
coordinate. If $Q$ is not a bijection and its output on uniform input is
$d$-wise uniform, the least possible excess in its ordered collision count is
\[
  \min_Q\bigl(\Coll(Q)-2^N\bigr)
    =2^{\max(\lceil N/k\rceil,\,d+1)}
    \qquad(k\ge1,\ 0\le d<N).
\]
A support bound gives the exponent $d+1$, and divisibility of Boolean
polynomial signed sums gives $\lceil N/k\rceil$. A ring with up to $k-1$
control bits on each edge attains both constraints simultaneously. We give
the quadratic case first, then prove the general result. For compressing
maps $Q\colon\F^n\to\F^m$ ($m\le n$) the minimum of $\Coll(Q)-2^{2n-m}$ is
$2^{\max(2\lceil n/k\rceil-m,\,\lceil n/k\rceil,\,d+1)}$, with a jump when $m$
drops below $\lceil n/k\rceil$. Lean~4 verifies both minima, including the
degree of the attaining maps.
\end{abstract}

\paragraph{Relation to prior work.}
Both lower bounds are classical: the support bound for functions whose
low-weight Fourier coefficients vanish~\cite{Delsarte73,MacWilliamsSloane}, and
McEliece-type divisibility of signed sums of degree-$k$
polynomials~\cite{McEliece72}. The contribution of this note is that the larger
of the two is always attained, by an explicit ring of gates, and the analogous
exact answer for compressing maps. We have not found these exact statements in
the literature; pointers are welcome.

\section{Introduction}

A map $Q\colon\F^N\to\F^N$ is a bijection exactly when its output on a uniform
input is uniform. Two natural weakenings of uniformity are \emph{$d$-wise
uniformity} (every $d$ output coordinates are jointly uniform) and \emph{low
collision count}. For arbitrary maps, these interact only through a
coding-theoretic support bound. For maps of algebraic degree two, a second
constraint appears: Walsh coefficients of quadratic functions are coarsely
quantised. Theorem~\ref{thm:main} gives the quadratic minimum. Theorem~\ref{thm:all-degrees} extends it to every positive degree bound, and Theorem~\ref{thm:rect} treats maps with fewer outputs than inputs.

\paragraph{Notation.}
For $c,y\in\F^N$ write $c\cdot y=\sum_i c_iy_i$ and $\wt(c)$ for the Hamming
weight. For $f\colon\F^N\to\mathbb{Z}$ let
$\hat f(c)=\sum_{y}f(y)(-1)^{c\cdot y}$ (unnormalised). For a map
$Q\colon\F^N\to\F^N$ let
\[
  \W_Q(c)=\sum_{x\in\F^N}(-1)^{c\cdot Q(x)},\qquad
  h(y)=|Q^{-1}(y)|,\qquad
  \Coll(Q)=\sum_y h(y)^2 .
\]
$Q$ is \emph{quadratic} if every coordinate is a polynomial of degree at most
two; then every component $c\cdot Q$ is a quadratic Boolean function.
We say $Q$ is \emph{silent through $d$} if $\W_Q(c)=0$ whenever
$1\le\wt(c)\le d$. By the XOR lemma this is equivalent to $Q(U)$ being $d$-wise
uniform.

\begin{definition}
For $N\ge1$ and $0\le d<N$ let $\mu(N,d)$ be the minimum of $\Coll(Q)-2^N$ over
quadratic maps $Q\colon\F^N\to\F^N$ that are silent through $d$ and are not
bijections.
\end{definition}

The restriction $d<N$ is necessary: silence through $N$ means $Q(U)$ is uniform,
so $Q$ is a bijection. Note $\Coll(Q)\ge 2^N$ with equality iff $Q$ is a
bijection, so $\Coll(Q)-2^N$ measures distance from bijectivity. Equivalently,
the collision probability of $Q(U)$ is $2^{-N}+(\Coll(Q)-2^N)\,2^{-2N}$.

\begin{theorem}\label{thm:main}
For all $N\ge 1$ and $0\le d<N$,
\[
  \mu(N,d)=2^{\max(\lceil N/2\rceil,\;d+1)} .
\]
\end{theorem}

\paragraph{Status.}
Both halves of the lower bound are classical (Lemmas~\ref{lem:support}
and~\ref{lem:divisibility}); we include short self-contained proofs. We were not
able to find the exact value $\mu(N,d)$, or the matching construction, in the
literature, and would be grateful for pointers.

\section{The deviation function}

Let $f(y)=h(y)-1$. Since $\sum_y h(y)=2^N$,
\begin{equation}\label{eq:excess}
  \sum_y f(y)=0, \qquad
  \sum_y f(y)^2=\Coll(Q)-2\cdot 2^N+2^N=\Coll(Q)-2^N .
\end{equation}
For $c\ne0$ we have $\sum_y(-1)^{c\cdot y}=0$, so
\begin{equation}\label{eq:walsh}
  \W_Q(c)=\sum_y h(y)(-1)^{c\cdot y}=\hat f(c)\quad(c\neq 0),\qquad \hat f(0)=0 .
\end{equation}
Thus $Q$ is a bijection iff $f\equiv 0$, and $Q$ is silent through $d$ iff
$\hat f(c)=0$ for every $c$ with $\wt(c)\le d$ (including $c=0$).

\section{Lower bounds}

\begin{lemma}[support bound]\label{lem:support}
Let $f\colon\F^n\to\mathbb{Z}$ be nonzero with $\hat f(c)=0$ for all
$\wt(c)<k$. Then $|\operatorname{supp} f|\ge 2^k$.
\end{lemma}

Multiplication by the full parity character sends each Fourier mask to its
complement. For $k\le n$ this gives a function of real multilinear degree at
most $n-k$ with the same support. In coding-theoretic terms this is a cube
support bound~\cite{Delsarte73,MacWilliamsSloane,ODonnell}. We prove it
directly below.

\begin{proof}
Induction on $n$. If $n=0$, then $f$ is a nonzero constant and $\hat f(0)=f\ne0$,
so $k=0$ and the claim holds. For $n\ge1$ and $k=0$ the claim is trivial. Let
$n,k\ge1$ and split on the last coordinate: $f_b(y')=f(y',b)$. Put
$s=f_0+f_1$ and $t=f_0-f_1$ on $\F^{n-1}$. Then
$\hat s(c')=\hat f(c',0)$ vanishes for $\wt(c')<k$, and
$\hat t(c')=\hat f(c',1)$ vanishes for $\wt(c')<k-1$.
If $s\neq0$, induction gives
$|\operatorname{supp} f|\ge|\operatorname{supp} s|\ge 2^k$.
If $s=0$, then $f_1=-f_0$ with $f_0\ne0$ and $t=2f_0$, so induction gives
$|\operatorname{supp} f|=2|\operatorname{supp} f_0|\ge 2\cdot2^{k-1}$.
\end{proof}

\begin{lemma}[quadratic divisibility]\label{lem:divisibility}
If $g\colon\F^n\to\F$ has degree at most two, then
$2^{\lceil n/2\rceil}$ divides $\sum_x(-1)^{g(x)}$.
\end{lemma}

This follows from Dickson's theorem: the Walsh values of $g$ lie in
$\{0,\pm2^{n-r/2}\}$ where $r$ is the (even) rank of its polar form
\cite[Ch.~15]{MacWilliamsSloane}, \cite{Carlet21}. Since $r\le 2\lfloor
n/2\rfloor$, the exponent is at least $\lceil n/2\rceil$. An elementary proof:

\begin{proof}
Induction on $n$; write $S(g)=\sum_x(-1)^{g(x)}$. If $g$ is affine then
$S(g)\in\{0,\pm2^n\}$. Otherwise some product $x_ix_j$ occurs; rename it $x_1x_2$
and write $g=x_1x_2+x_1A+x_2B+C$ with $A,B$ affine and $C$ quadratic in
$x_3,\dots,x_n$. Summing over $x_1$ gives $2[x_2=A]$, and then over $x_2$ gives
$2(-1)^{AB}$, so $S(g)=2\,S(C+AB)$, where $C+AB$ is quadratic in $n-2$ variables.
By induction $2^{1+\lceil (n-2)/2\rceil}=2^{\lceil n/2\rceil}$ divides $S(g)$.
\end{proof}

\begin{proof}[Proof of the lower bound in Theorem~\ref{thm:main}]
Let $Q$ be quadratic, silent through $d$, and not a bijection, and let $f$ be
its deviation. Then $f\ne0$.

\emph{Sparsity.} By~\eqref{eq:walsh}, $\hat f$ vanishes on all $c$ with
$\wt(c)<d+1$. Lemma~\ref{lem:support} gives
$|\operatorname{supp} f|\ge2^{d+1}$. Each nonzero integer $f(y)$ contributes at
least $1$ to~\eqref{eq:excess}, so $\Coll(Q)-2^N\ge 2^{d+1}$.

\emph{Quantisation.} Since $f\ne0$, Fourier inversion gives some $c$ with $\hat
f(c)\ne0$, and $c\ne0$ because $\hat f(0)=0$. By~\eqref{eq:walsh} and
Lemma~\ref{lem:divisibility} applied to $c\cdot Q$,
$|\hat f(c)|=|\W_Q(c)|\ge 2^{\lceil N/2\rceil}$. Since $f$ is integer-valued,
\[
  2^{\lceil N/2\rceil}\le|\hat f(c)|\le\sum_y|f(y)|\le\sum_y f(y)^2=\Coll(Q)-2^N.
  \qedhere
\]
\end{proof}

\begin{remark}
The sparsity bound uses nothing about degree, and the quantisation bound uses
nothing about silence. For arbitrary (non-quadratic) maps the sparsity bound
$2^{d+1}$ is already tight. For instance, the support-$2^{d+1}$ extremiser of
Lemma~\ref{lem:support} is realised by moving $2^d$ preimages. When $d=0$, a
quadratic non-bijection on $N\ge3$ bits cannot have a single double bucket and
a single empty bucket: such a map has excess $2<2^{\lceil N/2\rceil}$.
\end{remark}

\section{A matching construction}

Let $m=\max(\lceil N/2\rceil,d+1)$ and $s=N-m$. Then $1\le m\le N$ (as $d<N$)
and $s\le m$ (as $m\ge N/2$). Write inputs as $(t,w)$ with \emph{controls}
$t\in\F^s$ and \emph{wires} $w\in\F^m$ indexed by $\mathbb{Z}/m$. Put
$\tau_i=t_i$ for $i<s$ and $\tau_i=1$ for $s\le i<m$, and define
\[
  R(t,w)=(t,e),\qquad e_i=w_{i+1}+\tau_i\,w_i\quad(i\in\mathbb{Z}/m).
\]
Every coordinate has degree at most two.

\begin{lemma}\label{lem:ring}
For $t\ne\one$, the map $w\mapsto e$ is a bijection of $\F^m$. For $t=\one$
(that is, $\tau=\one$), every even-weight $e$ has exactly two preimages and every
odd-weight $e$ has none.
\end{lemma}

\begin{proof}
For fixed $t$ the map $w\mapsto e$ is linear. If $\tau_i=0$ for some $i$, then
$w_{i+1}=e_i$, and $w_{j+1}=e_j+\tau_jw_j$ determines every remaining wire in
turn around the ring. Hence the map is injective, and therefore bijective. If
$\tau=\one$, then $e_i=w_i+w_{i+1}$. The kernel is $\{0,\one\}$, and the image
lies in the even-weight words. By counting dimensions, the image is exactly the
even-weight words. (For $m=1$ this reads $e_0=0$.)
\end{proof}

\begin{proof}[Proof of the upper bound in Theorem~\ref{thm:main}]
By Lemma~\ref{lem:ring}, the deviation of $R$ is
\[ f(t,e)=[t=\one]\,(-1)^{\wt(e)} . \]
It is nonzero on exactly $2^m$ outputs, all
with $|f|=1$, so $\Coll(R)-2^N=2^m$ by~\eqref{eq:excess}. Its transform is
\[
  \hat f(a,b)=(-1)^{\wt(a)}\sum_{e\in\F^m}(-1)^{(b+\one)\cdot e}
  =(-1)^{\wt(a)}\,2^m\,[\,b=\one\,].
\]
Hence every nonconstant audible mask has $b=\one$ and weight at least
$m\ge d+1$. So $R$ is silent through $d$, is not a bijection, and attains
$2^{\max(\lceil N/2\rceil,d+1)}$.
\end{proof}

\begin{remark}
The extremiser has $2^{m-1}$ double buckets and $2^{m-1}$ empty buckets, all inside one
control slice. Its first audible mask is the all-wires mask, with coefficient
exactly $2^m$, so every inequality in the quantisation chain
$|\hat f(c)|\le\sum|f|\le\sum f^2$ is an equality; when $d+1\ge\lceil N/2\rceil$
the support bound is tight as well. For example,
$N=10$ and $d=4$ give $m=5$: $16$ double and $16$ empty buckets, every mask of
weight at most $4$ silent, and $\W_R(0,\one)=32$.
\end{remark}

\section{All positive degree bounds}\label{sec:all-degrees}

\begin{theorem}\label{thm:all-degrees}
For $k\ge1$ and $0\le d<N$, put $L=\max(\lceil N/k\rceil,d+1)$.
Among all square maps of coordinate degree at most $k$ that are silent
through $d$ and fail to be bijections, the minimum of $\Coll(Q)-2^N$ is $2^L$.
\end{theorem}

\begin{proof}
For the degree bound we use the signed-sum form of McEliece
divisibility~\cite{McEliece72}. Write a Boolean polynomial $g$ as a sum
modulo two of squarefree monomials $p$, each of degree at most $k$. Expand
$(-1)^{g(x)}=\prod_p(1-2p(x))$ over the integers. A term selecting $s$
monomials covers $u\le ks$ distinct variables. Its product equals one on
$2^{N-u}$ inputs, so it contributes $(-2)^s2^{N-u}$ to the signed sum.
Since $k(s+N-u)\ge u+(N-u)=N$, every term is divisible by
$2^{\lceil N/k\rceil}$. This applies to each $g=c\cdot Q$.
For a non-bijection some nonconstant coefficient is nonzero, and its
magnitude is at most $\sum_y f(y)^2$. Combining this degree bound with
Lemma~\ref{lem:support} gives $\Coll(Q)-2^N\ge2^L$.

For attainment, place $L$ wires on the ring and copy the $M=N-L$ controls
to the output. The inequality $N\le kL$ gives $M\le L(k-1)$. Partition
all $M$ controls into $L$ groups $S_i$ of size at most $k-1$, allowing
empty groups. Return the discrepancies
\[
  e_i=w_{i+1}+w_i\prod_{j\in S_i}t_j
  \qquad (i\in\mathbb Z/L\mathbb Z),
\]
with empty product one. Every coordinate has degree at most $k$. A zero
control makes its edge product zero, so Lemma~\ref{lem:ring} gives one
preimage for every discrepancy word on that slice. On the all-ones control
slice every edge product is one. Its $2^{L-1}$ even-parity outputs have two
preimages each and its $2^{L-1}$ odd-parity outputs have none. The deviation
is therefore $f(t,e)=[t=\one](-1)^{\wt(e)}$, with squared sum $2^L$.
Its transform is
\[
  \hat f(a,b)=(-1)^{\wt(a)}2^L[\,b=\one\,].
\]
The mask on the $L$ ring outputs alone reads $+2^L$, and every nonconstant
mask of smaller weight is silent. Since $d<L$, the ring attains the bound.
\end{proof}

This answers Question~2 of the quadratic draft. For instance, at $N=9$,
$k=3$, and $d=2$, each of the three edges takes two controls. Four double
buckets and four empty buckets give excess $8$. With the same $N,d$ the
quadratic minimum is $32$. The construction also covers $L=1$ (the
exceptional slice sends both wire values to zero) and $M=0$ (that slice
is the whole input cube).

\section{Compressing maps}\label{sec:rect}

Now let $Q\colon\F^n\to\F^m$ with $1\le m\le n$. Define $\W_Q$, $h$ and
$\Coll(Q)$ as before, with $y$ ranging over $\F^m$. A map whose output is
uniform has $\Coll(Q)=2^{2n-m}$, and this is the least possible value.

\begin{theorem}\label{thm:rect}
Let $1\le m\le n$, $k\ge1$ and $0\le d<m$, and put $\ell=\lceil n/k\rceil$.
Among all maps $Q\colon\F^n\to\F^m$ of coordinate degree at most $k$ that are
silent through $d$ and have $\Coll(Q)>2^{2n-m}$,
\[
  \min_Q\bigl(\Coll(Q)-2^{2n-m}\bigr)=2^{\max(2\ell-m,\;\ell,\;d+1)} .
\]
For $m\ge\ell$ the exponent is $\max(\ell,d+1)$. For $m<\ell$ it is $2\ell-m$,
which then exceeds both $\ell$ and $d+1$.
\end{theorem}

For $m=n$ this is Theorem~\ref{thm:all-degrees}, since $2\ell-n\le\ell$.

\begin{proof}[Proof of the lower bound]
Put $f(y)=h(y)-2^{n-m}$. As in Section~2,
$\sum_yf(y)^2=\Coll(Q)-2^{2n-m}$, $\hat f(0)=0$ and $\hat f(c)=\W_Q(c)$ for
$c\ne0$. Since $f\ne0$, some $c\ne0$ has $\W_Q(c)\ne0$, and $\W_Q(c)$ is a
multiple of $2^\ell$ by the divisibility step in the proof of
Theorem~\ref{thm:all-degrees}. This gives three bounds.
\begin{itemize}
\item $2^\ell\le|\hat f(c)|\le\sum_y|f(y)|\le\sum_yf(y)^2$, since $f$ is
  integer-valued.
\item Parseval: $2^m\sum_yf(y)^2=\sum_{c'}\hat f(c')^2\ge2^{2\ell}$, so the
  excess is at least $2^{2\ell-m}$.
\item Lemma~\ref{lem:support} on $\F^m$ gives
  $|\operatorname{supp}f|\ge2^{d+1}$, and each nonzero value has $f(y)^2\ge1$.
\end{itemize}
If $m<\ell$, then $2\ell-m>\ell>m>d$, which gives the second claim.
\end{proof}

\begin{proof}[Construction for $m\ge\ell$]
Let $L=\max(\ell,d+1)\le m$ and take the grouped ring from the proof of
Theorem~\ref{thm:all-degrees} on $n$ inputs: $L$ wires and $n-L$ controls,
with output $(t,e)$. Keep $e$ and a set $A$ of $m-L$ controls, and drop the
other $n-m$ controls, whose set we call $B$. Coordinate degrees do not grow.
Summing $h(t,e)=1+[t=\one](-1)^{\wt(e)}$ over $t_B$ gives
\[
  h'(t_A,e)=2^{n-m}+[t_A=\one]\,(-1)^{\wt(e)} .
\]
So the deviation $[t_A=\one](-1)^{\wt(e)}$ has the same shape as before. The
excess is $2^L$, the only nonzero coefficients are at $b=\one$ with
magnitude $2^L$, and the map is silent through $L-1\ge d$.
\end{proof}

\begin{proof}[Construction for $m<\ell$]
Take the grouped ring with $\ell$ wires and $n-\ell$ controls. This needs
$n-\ell\le\ell(k-1)$, which holds because $n\le k\ell$. Drop every control
and fold the surplus wires into the first output:
\[
  Q(x)=\pi(e),\qquad
  \pi(e)=(e_0+e_m+e_{m+1}+\dots+e_{\ell-1},\;e_1,\;\dots,\;e_{m-1}).
\]
The map $\pi$ is linear, so degrees do not grow. The histogram of $e$ is
$g(e)=2^{n-\ell}+(-1)^{\wt(e)}$. For $c\ne0$, the dual vector $\pi^{\top}c$
is nonzero, so
\[
  \W_Q(c)=\sum_e g(e)(-1)^{(\pi^{\top}c)\cdot e}
         =\sum_e(-1)^{(\pi^{\top}c+\one)\cdot e}
         =2^\ell\,[\pi^{\top}c=\one].
\]
Now $\pi^{\top}c=(c_0,c_1,\dots,c_{m-1},c_0,\dots,c_0)$ equals $\one$ only for
$c=\one$. The single audible mask has weight $m>d$, and by Parseval the excess
is $2^{2\ell}/2^m=2^{2\ell-m}$.
\end{proof}

For instance, take $n=12$ and $k=3$, so $\ell=4$. Every $m$ from $4$ to $12$
has minimum $2^{\max(4,d+1)}$, as in the square case. Below that the
minimum jumps: $2^5=32$ for $m=3$, $2^6=64$ for $m=2$ and $2^7=128$ for $m=1$.

\section{Formal verification}

Lean~4.15.0~\cite{Lean4} checks the proofs using its core and standard libraries
(no Mathlib). The definitions used in the two main statements are collected in
\texttt{ChallengeDefs.lean}, and \texttt{Challenge.lean} states
Theorems~\ref{thm:all-degrees} and~\ref{thm:rect} over them. A comparator check
confirms that \texttt{Solution.lean} proves exactly these statements over the
same definitions, so a reader only has to check those two short files.
For the quadratic case, \texttt{SilentRoom.lean} proves both lower bounds and
\texttt{RingCracks.lean} gives the ring histogram and spectrum.

For Theorem~\ref{thm:all-degrees}, \texttt{DegreeDivisibility.lean} proves
signed-sum divisibility from the monomial expansion and applies it to output
parities. \texttt{GroupedRing.lean} verifies the histogram, spectrum, and
coordinate degree of the same square map. \texttt{CubePermutation.lean}
relates positive excess to failure of bijectivity. In
\texttt{ExactCollisionMinimum.lean}, the theorem \texttt{exact\_minimum} gives
an attaining witness and a lower bound for every admissible map. Its axiom
report lists only \texttt{propext}, \texttt{Quot.sound}, and
\texttt{Classical.choice}.

For Theorem~\ref{thm:rect}, \texttt{RectangularCollisionMinimum.lean} proves
the three lower bounds, reusing the support bound and signed-sum divisibility.
It checks both constructions (the ring with dropped controls, and the folded
ring), including their coordinate degree. The theorem
\texttt{rect\_exact\_minimum} uses the same axioms.

The Lean files, the checking scripts and a script \texttt{verify.sh} that runs every
check are archived at \url{https://doi.org/10.5281/zenodo.24427416} and available at
\url{https://github.com/LizzoLizard/quiet-collisions-math}.

Independent enumeration checks every bucket, signed coefficient, and coordinate
degree for $495$ grouped rings through ten bits. Enumeration of all square maps
through three bits also confirms the minima, using output translations and
coordinate permutations to reduce the computation. Enumeration of every
eligible compressing map in $30$ cases with $n\le6$ confirms
Theorem~\ref{thm:rect}, in both regimes.


\section{Status of the original questions}

\begin{question}[rectangular maps, compression answered]
What is the minimum collision excess for $Q\colon\F^n\to\F^m$ with
$d$-wise uniform output, when $m<n$ (compression) or $m>n$ (stretch, as in
low-degree generators~\cite{MST06})? Theorem~\ref{thm:rect} answers the
compression case. For a stretch the output cannot be uniform, so the problem
needs a different baseline, and it remains open.
\end{question}

\begin{question}[higher degree, answered]
Is the degree-$k$ lower bound tight? Theorem~\ref{thm:all-degrees} gives the
exact minimum for every positive $k$.
\end{question}

\begin{question}[extremisers and robustness]
Is every extremal map affinely equivalent to the ring $R$? Does a robust version
hold, where low-weight coefficients are small rather than zero?
\end{question}

\section*{Acknowledgements}
I thank Claude (an AI model by Anthropic), which helped with the proofs, the
Lean formalization and the write-up, and Astra and Gemini Neon for their help.

{\small
\begin{thebibliography}{9}
\bibitem{Carlet21}
C.~Carlet, \emph{Boolean Functions for Cryptography and Coding Theory},
Cambridge University Press, 2021.

\bibitem{Delsarte73}
P.~Delsarte, An algebraic approach to the association schemes of coding theory,
\emph{Philips Research Reports Supplements} 10, 1973.

\bibitem{MacWilliamsSloane}
F.~J.~MacWilliams and N.~J.~A.~Sloane, \emph{The Theory of Error-Correcting
Codes}, North-Holland, 1977.

\bibitem{McEliece72}
R.~J.~McEliece, Weight congruences for $p$-ary cyclic codes,
\emph{Discrete Mathematics} 3, 1972.

\bibitem{MST06}
E.~Mossel, A.~Shpilka and L.~Trevisan, On $\epsilon$-biased generators in
$\mathrm{NC}^0$, \emph{Random Structures \& Algorithms} 29(1), 2006.

\bibitem{Lean4}
L.~de~Moura and S.~Ullrich, The Lean 4 theorem prover and programming language,
\emph{CADE-28}, 2021.

\bibitem{ODonnell}
R.~O'Donnell, \emph{Analysis of Boolean Functions}, Cambridge University Press,
2014.
\end{thebibliography}
}

\end{document}
